\section{Charges and anomalies}
\label{sec:charges}

The weak and colour algebras of a block are fixed by the projector; the abelian charges are not. This section asks how much the
algebra says about them, first inside the arena and then in a layer that uses nothing from the arena at all.

\subsection{A pairing obstruction}
\label{sec:pairing}

The family block of $\Pfam$ carries three abelian charges. The centre of $\keff(\Pfam)$ is spanned by the block images of
$X_a\in\{\tfrac12J^{(1)},\tfrac12J^{(3)},\tfrac12e^{(3)}_5e^{(3)}_6\}$, and their charges $q_a:=-J(2X_a)\Pfam$ are commuting symmetric
involutions of the block, with eigenvalues $\pm1$ and $\Tr q_a=0$. The cubic anomaly of $u=aq_1+bq_2+cq_3$ is the chiral trace
\[
\mathcal A(u^3)=\tfrac14\bigl[\Tr(u^3P_L)-\Tr(u^3P_R)\bigr],\qquad P_{L,R}=\Pfam\tfrac12(\id\mp\Gamma),
\]
per Weyl component, right-handed components entering with conjugated charge when every field is written left-handed. Since
$q_a^2=\Pfam$ and the linear chiral traces vanish, only the mixed term survives:
\[
\mathcal A(u^3)=6abc\,\mathcal A(q_1q_2q_3)=96\,abc ,
\]
while the plain trace $\Tr(u^3)$ vanishes for every $u$ and carries no information about anomalies. The anomaly-free locus is the union
of the three coordinate planes $abc=0$, and on it there is nothing to select.

\begin{theorem}[Conjugate pairing]
\label{thm:pairing}
Let a module carry three commuting involutions $q_a$ with eigenvalues $\pm1$ and $\Tr q_a=0$. For every $u=aq_1+bq_2+cq_3$ with
$abc=0$ the weight multiset of $u$ is symmetric, $\mathrm{mult}(w)=\mathrm{mult}(-w)$.
\end{theorem}

\begin{proof}
Take $c=0$; only the joint multiplicities $M_{s_1s_2}$ of $(q_1,q_2)$ matter. The trace conditions give
$M_{++}+M_{+-}-M_{-+}-M_{--}=0$ and $M_{++}-M_{+-}+M_{-+}-M_{--}=0$; adding and subtracting, $M_{++}=M_{--}$ and $M_{+-}=M_{-+}$. The
weight of $(s_1,s_2)$ is $as_1+bs_2$, and $(s_1,s_2)\mapsto(-s_1,-s_2)$ negates it while preserving multiplicities.
\end{proof}

The proof uses only the three linear traces, with no reference to rank, rung or realisation, so within its hypotheses no
enlargement of the block and no ascent of the tower repairs it. Charge pairing is all that is concluded: calling a full matter
representation vector-like would also require the chirality convention and the other gauge factors. And the hypotheses are not
decorative.

\begin{example}[Vanishing traces do not imply pairing]
\label{ex:multiset}
The multiset $\{(-3)^{\times5},1^{\times10},5^{\times1}\}$ has $\sum q=-15+10+5=0$ and $\sum q^3=-135+10+125=0$, yet it is not symmetric under
$q\mapsto-q$. It is realised by three commuting involutions with joint signs and multiplicities $(+,+,+){:}\,1$, $(-,-,+){:}\,10$,
$(+,-,-){:}\,5$: these have traces $(\Tr q_1,\Tr q_2,\Tr q_3)=(-4,-14,6)$, so the hypotheses of \cref{thm:pairing} fail, while
$Y=q_1+q_2+3q_3$ has spectrum $\{5^{\times1},1^{\times10},(-3)^{\times5}\}$ with $\Tr Y=\Tr Y^3=0$. Leaving the hypotheses therefore need not
violate the gravitational condition, which constrains $\Tr Y$ for the gauged combination and not $\Tr q_a$ for each basis charge
\cite{AlvarezGaumeWitten,Bilal}.
\end{example}

\subsection{Dressing, and a benchmark}
\label{sec:benchmark}

The escape from \cref{thm:pairing} is to leave its hypotheses. On the family block let
\[
t^1=-J\cdot\tfrac12e^{(2)}_5e^{(2)}_6,\qquad t^2=-J\cdot\tfrac12e^{(2)}_6e^{(2)}_4,\qquad t^3=-J\cdot\tfrac12e^{(2)}_4e^{(2)}_5
\]
be the weak charges, whose skew counterparts, the rotations of the spatial directions of constituent two, span the derived $\su2$ of
$\keff(\Pfam)$; let $\varepsilon=\pm1$ be an orientation sign, and put $t^a_{L,R}=P_{L,R}t^a$,
\[
B_0=\frac\varepsilon6\,\Gamma(q_1+q_2+q_3),\qquad Y=B_0+t^3_R,\qquad Q_{\rm em}=t^3_L+Y .
\]
The cell sums $\varepsilon\sum_aq_a$ are $\{-1,-1,-1,3\}$ on the $\Gamma=-1$ half and $\{-3,1,1,1\}$ on the other, which gives
$B_0|_{\Gamma=-1}=\{\tfrac16,\tfrac16,\tfrac16,-\tfrac12\}$ and, after adding $t^3_R=\pm\tfrac12$ on the right-handed cells,
$Y\in\{0,-1,\tfrac23,-\tfrac13\}$ there. Over the sixteen left-handed Weyl components $\Tr Y=\Tr Y^3=\Tr T_3^2Y=0$ with four doublets, so all
four anomaly conditions hold, and $\Tr T_3^2/\Tr Q^2=3/8$. This hypercharge is adopted, not derived.

\begin{proposition}[The escape]
\label{prop:escape}
\Cref{thm:pairing} constrains undressed elements of $\uu1^3$, which lie in the vector class. The charge $B_0$ carries a factor $\Gamma$ and is
axial, not an element of $\uu1^3$, so the cubic-trace formula does not apply to it. Once chirality-dressed endomorphisms are admitted,
which abelian charges can be gauged changes, and no-go results proved in the undressed grade-two setting do not carry across; a well-posed
search must say which Clifford grades are permitted alongside $Y^\dagger=Y$, $[Y,t_{La}]=0$ and $[Y,C_A]=0$.
\end{proposition}

Since $P_L+P_R=\id$, $Q_{\rm em}=t^3+B_0$: the chirality dressing cancels identically in the electric charge when the same axis is used in
both terms and is visible only in the split into $T_3$ and $Y$. The value $3/8$ is arithmetic on a charge table; it becomes $\sin^2\theta_W$
only after a common normalisation of the gauge kinetic terms and a matching scale, neither of which is supplied. Five choices enter the
benchmark: the projector, the dressing, which abelian direction is $Y$, which right-handed axis, and the overall sign.

\subsection{A model-independent layer}
\label{sec:modelindep}

The following uses nothing from the arena. It calibrates how much of the hypercharge problem is generic.

\begin{theorem}[The anomaly variety of one family]
\label{thm:variety}
For one family with $N_c$ colours and a right-handed neutrino, with left-handed hypercharges $q,l$ and right-handed $y_u,y_d,y_\nu,y_e$,
the three linear conditions give $l=-N_cq$, $y_d=2q-y_u$, $y_e=2l-y_\nu$. With $h_q=y_u-q$ and $h_\ell=y_\nu-l$ the cubic condition
factorises,
\[
\mathcal A_{Y^3}=6N_cq\,(h_\ell^2-h_q^2),
\]
using $2q^3-y_u^3-(2q-y_u)^3=-6qh_q^2$. The solution set is the union of three planes: $\Pi_A$ ($h_\ell=h_q$), $\Pi_B$ ($h_\ell=-h_q$) and
$\Pi_C$ ($q=l=0$). On $\Pi_A$ the four Yukawa relations $y_u=q+h$, $y_d=q-h$, $y_\nu=l+h$, $y_e=l-h$ hold, so they are consequences of anomaly
cancellation rather than assumptions; $\Pi_B$ is the image of $\Pi_A$ under the exchange $\nu_R\leftrightarrow e_R$. At $N_c=3$ and
$h=\tfrac12$, $\Pi_A$ is the line $Y=Y_{\rm SM}+\lambda(\BmL)$.
\end{theorem}

Anomaly cancellation does not supply the scalar, its hypercharge, or nonvanishing couplings. Quark Yukawa couplings alone do not exclude
$\Pi_C$ either: at $N_c=3$ the point $(q,l,y_u,y_d,y_\nu,y_e)=(0,0,\tfrac12,-\tfrac12,1,-1)$ with scalar hypercharge $\tfrac12$ cancels all four
anomalies and permits both quark Yukawas, yet $h_q=\tfrac12\neq h_\ell=1$. What excludes it is integrality: with $q=0$ the left-handed quarks
carry $\pm\tfrac12$ and a three-quark singlet has half-integral charge.

\begin{theorem}[Without $\nu_R$ hypercharge is quantised]
\label{thm:nonu}
Deleting $\nu_R$ gives $y_e=2l$ and $\mathcal A_{Y^3}=6N_cq(N_c^2q^2-h_q^2)$, so on the generic branch $h_q=\pm N_cq$ is locked and
hypercharge is fixed up to scale and relabelling: at $N_c=3$, $q=\tfrac16$ gives the Standard-Model values with no residual $\BmL$ direction.
\end{theorem}

This is the classical charge-quantisation result \cite{FootLewVolkas,GengMarshak}. Its bearing here is that the difficulty of selecting
hypercharge on a family block exists because the block carries sixteen Weyl components rather than fifteen: the sterile neutrino is both
the block's virtue and the source of the ambiguity.

\begin{proposition}[Integrality, and what it does not fix]
\label{prop:integrality}
On $\Pi_A$ with $Q_{\rm em}=T^3_L+Y$, the leptons carry $Q_\nu=-\lambda$ and $Q_e=-1-\lambda$, so integrality of the charged lepton forces
$\lambda\in\Z$, and colour triality then gives integral charge to every colour singlet. It does not fix the fractional part of the quark
charges, since $Q_u=\tfrac23+\lambda/3$, and neutrality does not close the argument for composites: $Q(ddd)=\lambda-1$ and $Q(uuu)=\lambda+2$, so
neutral colour-singlet operators of odd fermion parity exist at $\lambda=1$ and $\lambda=-2$. A neutral \emph{elementary} colour-singlet
fermion exists if and only if $\lambda\in\{0,-1\}$, and those two values are exchanged by $\nu\leftrightarrow e$.
\end{proposition}

Three other routes fail to single out the Standard Model. The kernel of the representation is not a selection principle: its order
depends on $\lambda$ through the integer normalisation of the weights. Nor is a $\Z_6$ kernel evidence for three colours: with a neutral
$\nu_R$ the kernel is $\Z_{2N_c}$ for odd $N_c$ and $\Z_{N_c}$ for even $N_c$. And integrality does not select three either: on the standard
branch every $N_c$-quark colour singlet has integral charge exactly for odd $N_c$, and the number $N_c+1$ of weak doublets per family, which
the global $SU(2)$ anomaly requires to be even \cite{Witten82}, is even exactly for odd $N_c$. All three facts hold at every odd $N_c$. With
$N_g$ generations the standard branch has $2N_g$ parameters, and the anomaly-free directions include $L_e-L_\mu$ and $L_\mu-L_\tau$
\cite{HeJoshiLewVolkas}, so family universality is an input that grows with $N_g$. The global structure of the gauge group is a separate
datum \cite{Tong,GaiottoKapustinSeibergWillett}.

\subsection{Charge at one constituent}

Complexifying, the spinor module of one constituent is $\Lambda^\bullet\C^3$, and its even half is spanned by
$|0\rangle,|12\rangle,|13\rangle,|23\rangle$. The diagonal Fock charge $N_1-N_2$ has spectrum $\{0,0,+1,-1\}$ there, the charge skeleton of
$(\nu,\bar\nu,e^+,e^-)$.

\begin{theorem}[No unique charge at one constituent]
\label{thm:fock}
The condition $\spec Q=\{0,0,+1,-1\}$ on the even half, for $Q=\sum_iq_iN_i+c$, has exactly twelve solutions, forming two orbits of
$S_3\times\{\pm1\}$: $Q=N_i-N_j$ with $c=0$, and $Q=N_i+N_j-1$ with $c=-1$. The antiparticle condition $CQC^{-1}=-Q$, with $C$ the Hodge
conjugation, holds on both orbits and does not separate them.
\end{theorem}

\begin{proof}
Summing the four values gives $4c+2\sum_iq_i=0$. If the neutral pair is $\{|0\rangle,|ij\rangle\}$ then $c=0$ and $q_i+q_j=0$, which forces
the third coefficient to vanish; if it is $\{|ij\rangle,|ik\rangle\}$ then $q_j=q_k$, $c=-q_j$ and $q_i=0$.
\end{proof}

The further step, $Q|0\rangle=0$ and hence $c=0$, is the separate assumption that the state identified with the neutrino is the Fock vacuum.
Choosing between the orbits is an identification, not a derivation, and it has the same shape as the hypercharge problem two rungs up,
differing in size and not in kind.
